\chapter{绪论}\label{chap:Intro}

\section{研究背景及意义}\label{sec:Intro_background_and_significance}

% 研究背景

能源背景-电动汽车-电动汽车关键部件的PHM-PHM的前世今生，为什么研究智能化PHM

% 研究意义，本质是要解决的问题为什么困难

智能化PHM面临哪些挑战

\section{国内外研究现状}\label{sec:Intro_literature_review}

In this section, we review methodologies relevant to uncertainty-aware prognosis, categorizing them into two paradigms: physics-based models and data-driven models.

In this section, we review methods related to robust battery degradation modeling, including methods based on (1) pretraining-finetuning paradigm, (2) domain adaptation paradigm, and (3) domain generalization paradigm.
% (1) pretraining-finetuning requiring full target-domain labels,
% (2) domain adaptation aligning feature distributions with unlabeled target data,
% (3) domain generalization operating without target-domain supervision.

% In this section, we review the related methods including CDLMs and LLMs for mitigating statistical bias, CDLMs and LLMs for rejecting undesired data, and LLMs for detecting sequence-level hallucinations.
In this section, we review the related methods including (1) CDLMs and LLMs for mitigating statistical bias, (2) CDLMs and LLMs for rejecting undesired inputs, and (3) LLMs for detecting sequence-level hallucinations.

\subsection{计及不确定性的健康状态估计方法研究现状}\label{subsec:SOH_estimation}

（1）Uncertainty-aware Physics-based Prognostic Models % \label{subsec:physics-based_models}

Physics-based prognostic models derive uncertainty-aware predictions by integrating degradation mechanisms with filtering techniques for state estimation and parameter updating. The employed filters—predominantly Kalman filters and particle filters—leverage inherent Bayesian frameworks to quantify uncertainty through probabilistic inference over system states and noise distributions\cite{chen2003bayesian}, enabling dynamic confidence bounds on predictive results.
% 2016
% Bressel等人\cite{7386653} propose an extended Kalman filter-based algorithm integrated with the inverse first-order reliability method for SOH and RUL prediction, along with the corresponding confidence intervals.
文献\cite{7386653}中，Bressel~等人提出了一种基于观测器的预测方法，所提出的方法通过扩展卡尔曼滤波器（Extended Kalman Filter, EKF）建模设备的健康状态和退化动态，基于逆一阶可靠度方法（Inverse First-Order Reliability Method, IFORM）外推健康状态轨迹直至触发预设的失效阈值，从而实现对变功率负载工况下的质子交换膜燃料电池的剩余使用寿命预测，并给出90\%置信区间作为对预测结果的不确定性的量化。
% 2021
% Kim等人\cite{9032349} propose a prognostic model utilizing a dynamic particle filter which quantifies the uncertainty associated with the battery's end-of-life and adaptively adjusts the parameters of the degradation model.
文献\cite{9032349}中，Kim~等人结合异常检测技术和基于粒子滤波（Particle Filter, PF）的预测方法提出了一种平移核粒子滤波（Shifting Kernel Particle Filter, SKPF）算法，解决了特定循环阶段带有异常退化轨迹的电池的突变特征辨识问题，实现了兼具精度和稳健性的剩余使用寿命预测，所提出方法同样给出预测结果的置信区间。
% 2021
% Haque and Choi \cite{9113648} propose a prognostic model utilizing a sparse kernel ridge regression-assisted particle filter and quantifie abrupt changes in degradation trajectories by accounting for noise uncertainty.
晶体管的可靠应用依赖剩余使用寿命预测，晶体管性能退化具有不确定性、突变性和样本贫化性，这与电池性能退化在数据表现上一致，文献\cite{9113648}中，Haque~和~Choi~提出了一种稀疏核岭回归辅助的粒子滤波（Sparse Kernel Ridge Regression Assisted Particle Filter, SKRR-PF）算法，通过在后验概率密度函数估计的重采样环节引入稀疏核岭回归机制，实现了对晶体管器件导通电阻轨迹突变的识别和剩余使用寿命的预测。
% 2022
% Ahwiadi and Wang \cite{AHWIADI2022110817} propose an enhanced particle filter model based on evolutionary fuzzy predictors to mitigate sample degeneracy and impoverishment issues while achieving high-precision remaining useful life interval estimation.
文献\cite{AHWIADI2022110817}中，Ahwiadi~和~Wang~提出了一种增强型粒子滤波（Enhanced Particle Filter, EPF）框架，通过融合增强型粒子机制和演化模糊预测器（Evolving Fuzzy Predictor, EFP），克服基于传统~PF~方法进行电池系统状态估计和寿命预测时由权值退化和粒子多样性丧失导致的建模不确定性，从而实现准确且鲁棒的电池系统健康状态估计和剩余寿命预测。
% 2023
% Rezaei等人\cite{REZAEI2023107883} develope a fuzzy robust two-stage unscented Kalman filter to estimate cycle capacity and quantify uncertainty in a phased manner using the noise covariance matrix.
文献\cite{REZAEI2023107883}中，Rezaei~等人提出了一种模糊鲁棒两阶段无迹卡尔曼滤波器（Fuzzy Robust Two-Stage Unscented Kalman Filter, FRTSUKF），实现对缺乏先验统计特性的电池模型进行不确定性建模，其在第一阶段基于模糊推理系统剔除测量噪声，初步估计电池模型状态变量，在第二阶段进行不确定性量化，并基于量化结果补偿电池性能退化向状态变量引入的不确定性。
% 2024
% Cui等人\cite{10385053} design a prognostic model based on dual unscented particle filters to address uncertainties in dynamic degradation trajectories and achieve robust prediction.
旋转机械设备维护过程同样存在状态估计和寿命预测需求，且设备性能退化行为与电池系统一致。文献\cite{10385053}中，Cui~等人提出了一种双无迹粒子滤波器（Dual Unscented Particle Filter, DUPF），通过集成双流异构模型并设计全局融合策略，实现跨时间尺度轴承退化特征，并通过显式量化动态贝叶斯最大失效概率提升了预测分布的统计可靠性。
% 2024
% Wang等人\cite{10672556} propose a zonotopic Kalman filter-based method with physics constraints to achieve online capacity interval estimation.
文献\cite{10672556}中，为了实现电池在线容量估计，Wang~等人提出了一种物理约束中心对称多胞形卡尔曼滤波器（Physics-Constrained Zonotopic Kalman Filter, PCZKF），其中所提出的滤波器实现电池模型的参数及其区间估计的联合更新，所提出的物理约束则用于应对电池容量再生挑战，所提出的方法相比传统基于~KF~的方法能够输出更加精确容量预测轨迹和更加收敛的容量置信区间。

（2）Uncertainty-aware Data-driven Prognostic Models % \label{subsec:data-driven_models}

Uncertainty-aware data-driven prognostic models can be classified into statistical learning-based methods \cite{van2021bayesian, gruber2025sourcesuncertaintysupervisedmachine} and deep learning-based methods \cite{ABDAR2021243}. Statistical learning approaches rely on explicit probabilistic models for analytical or approximate inference. In contrast, deep learning methods address uncertainty by leveraging the statistical properties of randomized parameter spaces or output distributions.
% 
相关向量回归（Relevance Vector Regression, RVR）和高斯过程回归（Gaussian Process Regression, GPR）是两种典型的具有不确定性量化能力的统计学习方法，被广泛应用在状态估计和寿命预测问题中。
% Wang等人\cite{rvr_rul_9552552} introduce a prognostic method based on relevance vector regression (RVR), while Richardson等人\cite{gpr_soh_8263147} employ Gaussian process regression (GPR) for SOH estimation.
文献\cite{rvr_rul_9552552}中，Wang~等人针对传统~RVR~方法局限于一维退化过程和单变量观测数据的问题，提出了一种基于动态多元相关向量回归（Dynamic Multivariate Relevance Vector Regression, DMRVR）的退化轨迹跟踪和剩余寿命预测框架，所提出方法考虑多元物理环境因素，构建多步回归架构以表征退化动力学，并依据RVR中间结果的矩阵分布基于加速梯度算法优化求解过程。
% 
文献\cite{gpr_soh_8263147}中，Richardson~等人面向恒流充放电工况提出了一种基于高斯过程回归的原位容量数据驱动诊断方法（Gaussian Process Regression for \textit{In Situ} Capacity Estimation, GP-ICE），其在特定的电压窗口内，仅需提取短至十秒的连续恒流工况片段即可实现高精度电池健康状态估计。所提出的方法是一种贝叶斯非参数化回归方法，能够输出预测结果的不确定性作为对数据复杂性的表征。
% 
% Zhang等人\cite{bbb_rul_9738999} develop a Bayesian mixture neural network (BMNN) to capture the distribution of prediction results. 
文献\cite{bbb_rul_9738999}中，Zhang~等人提出了一种贝叶斯混合神经网络（Bayesian Mixture Neural Network, BMNN）以应对因内部反应非线性和高混杂性导致的剩余寿命预测不确定性大的问题。所提出网络引入准格雷姆角场（Quasi-Gramian Angular Field, Q-GAF）进行信号预处理以增强先验分布辨识能力，并融合贝叶斯卷积神经网络（Bayesian Convolutional Neural Network, BCNN）和贝叶斯长短期记忆网络（Bayesian Long Short-Term Memory Network, B-LSTM），由变分贝叶斯推断层输出预测结果的后验概率分布。
%
% Li等人\cite{capsnet_rul_li2024sensor} propose a dropout-based remaining useful life (RUL) prediction framework ，命名为SACN.
为了解决考虑不确定性的预测方法存在的模型合理性和可解释性欠缺的问题，文献\cite{capsnet_rul_li2024sensor}中，Li~等人提出了一种考虑传感器的胶囊神经网络（Sensor-Aware Capsule Neural Network, SACN），这一工作提出了基于胶囊神经元激活水平的蒙特卡洛随机失活机制（Capsule-Activation-Based Monte-Carlo Dropout, CA-MCD），并开展了考虑跨传感器相关特征分布的不确定性归因研究。
%
% Zhang等人\cite{ensemble2_soh_10102497} combine a multi-model fusion method and a deep ensemble neural network to quantify SOH estimation uncertainty，命名为DeNN.
文献\cite{ensemble2_soh_10102497}中，Zhang~等人针对包含复杂环境噪声的动力电池循环信号，提出一种面向任意使用工况的数据驱动多模型融合健康状态评估方法，命名为DeNN。该方法利用模型集成（Model Ensemble, ME）技术并利用~KF~实现多个模型的输出的融合，既实现了多样负载工况下准确的健康状态估计，又实现了不确定性量化。
%
% Zhou等人\cite{mcdropout_autoencoder_10412192} design a Bayesian deep autoencoder (BDAE) with Monte Carlo dropout (MCD) and variational inference to measure the uncertainty in SOH estimates.
文献\cite{mcdropout_autoencoder_10412192}中，Zhou~等人提出了一种基于MCD的贝叶斯深度自编码器（Bayesian Deep Autoencoder, BDAE）以克服锂离子电池复杂且高度时变的实际运行工况带来的健康诊断挑战。该编码器首先进行自动化退化特征提取，继而开展优选输入数据筛选并构建回归模型，最后对预测结果关联的不确定性实施定量表征。
%
% Nemani等人\cite{nemani2023uncertainty} propose a strategy focusing on the spectral-normalized neural Gaussian process (SNGP) for deterministic uncertainty quantification.
基于蒙特卡洛随机失活机制（Monte-Carlo Dropout, MCD）的确定性-随机模型变换方法需要多次后采样，不利于深度学习模型的实时应用。文献\cite{nemani2023uncertainty}中，Nemani~等人综述了BDL的主要分支，并提出了一种基于谱归一化神经高斯过程（Spectral-Normalized Neural Gaussian Process, SNGP）的考虑不确定性的寿命预测框架。该框架将铺归一化应用于模型隐藏的残差层并使用高斯过程层替换全连接层，从而实现不确定性量化，并具有与单一确定性网络相当的推理延迟水平。
%
% Lin and Li\cite{9729404} introduce a Bayesian neural networks (BNN) framework utilizing Concrete Dropout for RUL prediction, in which an iterative calibration is integrated by combining isotonic regression with standard deviation (STD) scaling.
文献\cite{9729404}中，Lin~和~Li~面向可靠性决策对模型预测结果全局置信水平准确量化的需求，提出了一种基于贝叶斯深度学习（Bayesian Deep Learning, BDL）的剩余使用寿命预测框架。所提出的框架首先给出认知不确定性和偶然不确定性的数学表征和解耦方法，继而建立基于长短期记忆神经网络（Long-Short Term Memory Network, LSTM）的贝叶斯模型，该框架最后提出了一种迭代式校准策略，基于保序回归和标准差缩放（Standard Deviation Scaling, SDS）技术实现对认知不确定性、偶尔不确定性和全局预测不确定性的联合校准。
%
% Ke等人\cite{10217044} develop a MCD-based BNN model for LIBs' knee capacity prediction and also used the STD scaling method for uncertainty calibration.
文献\cite{10217044}中，Ke~等人提出了一种基于贝叶斯神经网络（Bayesian Neural Network, BNN）的拐点-拐点容量早期预测框架。该框架引入一种基于容量矩阵滑动窗口的退化特征提取机制，继而基于BNN和MCD实现锂离子电池缓慢老化-加速老化拐点的容量预测及对应的不确定性量化，最后基于SDS技术标定所量化的不确定性。
%
% Jiang等人\cite{10391265} present a neural network based on Bayes-by-Backprop strategy with uncertainty calibration to estimate SOH of second-life batteries.
储能系统是锂离子电池的一个重要应用领域，其使用电动汽车退役锂电池，也即梯次利用锂电池。在这个应用场景中，大规模处于不同生命周期的锂电池通过串联或并联在一起提供均衡稳定电能。储能系统中电池初始健康状态不尽相同，且不同局部的锂电池经历不同的充放电工况，导致了数据驱动的健康状态估计模型的评估不确定性。
针对此，文献\cite{10391265}中，Jiang~等人提出了一种梯次利用锂离子电池健康状态评估方法，该方法设计了一种联合特征提取策略以从电池循环序列中结构包含工况无关退化机理的健康状态表征信息，随后，基于具备不确定性校准能力的BNN生成全局健康状态估计结果，并定量表征模型评估过程的不确定性。

（3）研究的局限性

Though advancements, there remains several open issues in delivering prognostic results that are both highly accurate and free from bias.
% 
First, the physics-based models mainly focus on noise-caused uncertainties, leading to insufficient representation of model confidence and inconsistent outcomes. Conversely, those in data-driven models lack granularity and structural adaptability, impeding the effective extraction of degradation patterns.
% 
Second, the existing methods that focus on model uncertainty in predicting SOH often neglect spatial-temporal uncertainties in the data. This can lead to the use of irrelevant or incorrect data segments in degradation modeling, consequently diminishing the prognostic accuracy.
% 
Third, due to their approximate nature, uncertainty-aware prognostic models may inadequately represent the actual distribution of cyclic data. They exhibit small uncertainty even when their estimates substantially deviate from actual values, thus failing to differentiate between ID and OOD samples.
% 
To address the aforementioned issues, we propose HSSTT, a method for trustworthy SOH estimation that leverages uncertainty modeling and incorporates a calibration process to enhance both attention aggregation and the extraction of spatial-temporal degradation dependencies.

\subsection{面向分布漂移的退化建模方法研究现状}\label{subsec:degradation_modeling}

（1）{Pretraining-finetuning-based Degradation Modeling}

% In \cite{LI2021116410}, Li等人 propose a lightweight and generalizable model for battery SOH estimation that integrates the pretraining-finetuning paradigm with pruned convolutional neural network.
In \cite{9343713}, Che等人 introduce an online self-calibrating model for RUL prediction using Gaussian process regression, long short-term memory neural network (LSTM), and a pretraining-finetuning strategy.
In \cite{LIU2025115347}, Liu等人 fine-tune a bidirectional LSTM-Mamba2 model with differential voltage features to achieve SOH estimation across batteries with different operating conditions and electrode materials.
In \cite{D2EE01676A}, Ma等人 propose a real-time personalized battery degradation modeling framework named CRNN, achieving joint prediction of SOH and RUL across operating conditions.
In \cite{wang2024physics}, Wang等人 propose a physics-informed neural network (PINN) for battery degradation modeling and achieve cross-dataset battery health state prediction using the pretraining-finetuning paradigm.

（2）{Domain Adaptation-based Degradation Modeling}

In \cite{9788003}, Ma等人 propose a method named ConvMMD that aligns source and target domains by minimizing the maximum mean discrepancy of lower-order moments between them.
In \cite{HAN2022230823}, Han等人 present a model requiring only a few target samples, achieved by incorporating a domain adaptation layer into LSTM, known as LSTM-DA.
In \cite{9560040}, Ye等人 design a method named DDAN that realizes unsupervised degradation feature alignment via adversarial training across domains.
In \cite{10029904}, Lu等人 develop a two-stage prognostic framework named DKTL, which combines degradation mechanisms and achieves knowledge transfer of degradation patterns from synthetic data to real data.
In \cite{WANG2023108897}, Wang等人 develope a model with a domain invariant-specific feature decoupling strategy and a degradation trend preservation strategy to align degradation patterns, known as DR-Net.

（3）{Domain Generalization-based Degradation Modeling}

Domain generalization paradigm involves three main categories, namely data manipulation, representation learning, and specific learning strategy \cite{9782500}.
% In \cite{CHEN2024234696}, Chen等人 develop a dual-branch feature extractor to process domain samples with mixed and separated sources. By leveraging a domain discriminator and a health state estimator, their approach enables domain-invariant degradation feature extraction, culminating in a representation learning-based degradation pattern generalization method termed DGSOH.
In \cite{CHEN2024234696}, Chen等人 develop a dual-branch network, a domain discriminator and a health state estimator to learn domain-invariant degradation feature extraction, culminating in a representation learning-based degradation pattern generalization method termed DGSOH.
In \cite{TAN2024103725}, Tan等人 design a meta learning-based domain generalization method named MAGNet. It couples a two-phase meta-optimization strategy with RUL-driven health prediction, enforcing domain-invariant representation learning through task-aligned gradient synchronization. % promoting the learning of cross-domain consistent representations.
In \cite{liu2024itransformer}, Liu等人 propose a method designed to obtain robust representations against OOD temporal patterns through the inclusion of a transposed attention module and a feed-forward module, known as iTransformer.
In \cite{nie2023a}, Nie等人 present a sequence modeling framework resistant to pattern drift known as PatchTST.
% The model employs subsequence segmentation and channel-independent modeling to enhance its adaptability to dynamic temporal patterns.
In \cite{chen2023tsmixer}, Chen等人 develop a linear model capable of efficiently extracting spatio-temporal dependencies without relying on task-specific prior knowledge, known as TSMixer.

（4）研究的局限性

Due to dependency on target domain data, pretraining-finetuning and domain adaptation paradigms face deployment constraints in scenarios with dynamic operating conditions and partial observability.
% Representation learning and meta-learning strategies eliminate the need for target domain samples but often rely on precise domain labels or complex pretext tasks, limiting their effectiveness in intricate transportation environments.
While the representation learning-based and meta learning-based methods enable target-free modeling, their dependence on degradation priors, including auxiliary metadata or pretext tasks, limits scalability and generalization performance.
% 
% These limitations inspire the exploration of methods driven by the relatively understudied data manipulation strategy, particularly the feature manipulation strategy.
Though promising, data manipulation-based method faces open issues in robust battery prognosis, including scarce spatiotemporal degradation information, statistics uncertainty in feature spaces due to domain shifts, and the long-tail degradation distributions across multi-source domains.
To address these issues, we propose RODTNet, a novel model for OOD degradation pattern generalization. It generates continuous degradation patterns within a low-complexity latent feature space through heuristic feature reconstruction and probabilistic feature statistics generation. Furthermore, we propose RODTNet-LT, an enhanced variant of RODTNet, which incorporates statistical feature correction to specifically address the long-tail distribution challenge.

\subsection{幻觉鲁棒的故障诊断方法研究现状}\label{subsec:fault_diagnosis}

（1）{CDLMs and LLMs for Mitigating Statistical Bias}

In \cite{Peng_Liu_Du_Gao_Wang_2025}, H. Peng等人 design a cross-dataset diagnostic LLM, termed BearLLM.
In \cite{VGCDM}, X. Liu等人 introduce pulse-voltage embeddings into the diffusion process for controllable vibration synthesis.
In \cite{DCGAN}, K. Zhou等人 develop a balanced generator-discriminator architecture to diagnose gear defects, termed DCGAN.
In \cite{WOS:000688310200005}, T. Han等人 propose a hybrid network that combines class constraints with adversarial domain regularization to capture transferable features.
In \cite{10478558}, M. Vu等人 propose a dual-branch diagnostic framework that enhances fault classification accuracy in sparse-label scenarios.
In \cite{QSCGAN}, W. Wan等人 embed self-attention and spectral normalization within generator-discriminator architectures to augment rocket bearing data, termed QSCGAN.
In \cite{VQVAE}, Q. Zheng等人 combine degradation stage partitioning and vector quantized augmentation to enrich degradation features, termed VQVAE.

（2）{CDLMs and LLMs for Rejecting Undesired Inputs}

%MDCC=DSRC+马氏距离
%WDSN=DANN+马氏距离
In \cite{ZHAO2022108672}, C. Zhao等人 employ a triplet loss and a boundary adaptation method to learn compact features and optimal class boundaries.
In \cite{10262196}, Z. Zhu等人 develop an adversarially trained open-set classifier to identify unknown classes and out-of-domain samples.
In \cite{10214410}, J. Long等人 develop a domain-generalized unseen fault diagnosis method by fusing a meta-learning block with the Mahalanobis metric, termed MLMM.
In \cite{ZHANG2023109518}, X. Zhang等人 introduce a domain-adaptive approach that disentangles private and shared features to diagnose unseen faults, termed WDSN.
In \cite{10382502}, B. Lu等人 design a contrastive embedding-based diagnostic strategy that jointly preserves domain-specific and domain-invariant features, termed MDCC.
In \cite{20242516270368}, B. Liu等人 benchmark LLaMA-series LLMs for undesired input detection and design a cosine distance-based detector using isotropic embeddings, termed LLM-OOD.

（3）{LLMs for Detecting Sequence-Level Hallucinations}

%请检查LLM-check中具体用到的特征，确保与实验一致。
%请检查\cite{azaria-mitchell-2023-internal}是否是秒杀级对比算法 <- 这篇论文没开源
%请检查这篇论文Detecting Hallucinations in Large Language Model Generation: A Token Probability Approach，参考文献中有好多利用中间层信息进行幻觉检测的文献。
%请检查Unsupervised Real-Time Hallucination Detection based on the Internal States of Large Language Models是否是秒杀级对比算法 <- 使用LLM-Check-HiddenScore实现
In \cite{20251218102787}, G. Sriramanan等人 design a compute-efficient detector that leverages the model's internal attention maps, termed LLM-Check.
In \cite{20243917091282}, J. Duan等人 develop a hallucination quantification approach by reweighting token- and sequence-level relevance.
In \cite{20240715556058}, T. Zhang等人 introduce a retrieval-free detector that quantifies hallucinations through token-level uncertainty.
In \cite{selfcheckgpt}, P. Manakul等人 propose SelfCheckGPT, which assesses hallucinations by sampling multiple black-box LLM responses and measuring their mutual consistency.
% In \cite{azaria-mitchell-2023-internal}, A. Azaria等人 design a hallucination classifier using LLM hidden activations.
In \cite{su-etal-2024-unsupervised}, W. Su等人 propose a real-time and annotation-free hallucination detection framework using the internal states of LLMs, termed MIND.
In \cite{10.1007/978-3-031-86623-4_13}, E. Quevedo等人 propose a supervised approach that utilizes an auxiliary detector trained on resource-efficient numerical features, termed SHalluDetect.
% In \cite{20242516270368}, B. Liu等人 highlight the isotropic nature of LLMs and design a cosine distance-based hallucination detector, termed LLM-OOD.
In \cite{li-etal-2023-halueval}, J. Li等人 develop a benchmark via sampling-then-filtering and annotation for hallucination analysis.

（4）研究的局限性

However, several open issues remain in improving the performance of diagnostic LLMs.
First, existing models overlook the compound statistical bias caused by inter-dataset discrepancies and class imbalance, often mistaking domain-specific features for stochastic noise.
Simple perturbation-based augmentation further aggravates this issue by generating strongly correlated samples that do not effectively enrich minority classes.
Second, CDLMs and LLMs tend to underestimate the predictive uncertainty associated with undesired inputs that share high-level semantics with the training data, thereby degrading diagnostic accuracy.
Third, due to semantic redundancy in generated text and the model's overconfidence, existing detectors frequently fail to identify hallucinated content.
To tackle these issues, we propose DA-TMLLM, which detects and mitigates LLM hallucinations by effectively augmenting the training data and accurately quantifying the uncertainty of generated answers.

\section{研究目标与关键研究挑战}\label{sec:Intro_research_objectives_and_challenges}

围绕研究目标，本文将研究中心落在三个相互关联的挑战上，以兼顾

关键挑战如下：

（1）

（2）

（3）

\section{研究内容与技术路线}\label{sec:Intro_research_content_and_approach}

基于此研究路线，本文的具体研究内容分为以下三个部分：

（1）

该方法旨在

（2）

该方法旨在

（3）

该方法旨在

\section{论文组织结构}\label{sec:Intro_thesis_structure}

本文一共分为六个章节，剩余章节的内容安排如下：

第二章

通过多个数据集上的对比实验验证了所提出方法在健康状态估计精度和不确定性量化质量方面的有效性和先进性。

第三章

大量定性和定量的实验结果验证了所提出方法在退化模式泛化问题上的有效性和先进性。

第四章

所提出的方法在故障诊断精度和幻觉检测准确性方面稳定优于现有方法。

第五章对全文进行总结与展望。
